\section{Compactness with a common lower barrier}

A common lower barrier makes all finite and infinite rooftops genuine. An error integral against each approximating measure bounds the rooftop error by a summable tail. Weighted lower semicontinuity passes this bound to the limit, and singular domination identifies the lower envelope with the $L^1$ limit. Sandwiching by upper envelopes then gives capacity convergence; suitable domination also controls the measure tails.

\Needspace{16\baselineskip}
\begin{theorem}\label{thm:common-barrier-criterion}\index{Barrier!common lower bound}
Let $\theta\geq0$ be big. Suppose $\varphi\in\PSH(X,\theta)$ admits a genuine $P_{\delta\theta}(\varphi)$ for some $0<\delta<1$. Let $u_j\in\PSH(X,\theta)$ satisfy, for a fixed $C$,
\[
 \varphi-C\leq u_j\leq\varphi+C,\qquad u_j\longrightarrow u\text{ in }L^1(X).
\]
If
\begin{equation}\label{eq:common-barrier-self-error}
 I_j:=\int_X|\mathrm{e}^{u_j}-\mathrm{e}^u|\,
               \mathrm{MA}^{\mathrm{ext}}_\theta(u_j)\longrightarrow0,
\end{equation}
then $u_j\to u$ in capacity.
\end{theorem}
\begin{proof}
The canonical limit lies between $\varphi-C$ and $\varphi+C$. It suffices to prove capacity convergence after further extraction from an arbitrary subsequence. Choose an extraction with $I_j\leq2^{-j}$ and almost-everywhere convergence for smooth volume. For $N\geq j$ define
\[
 p_{j,N}=P_\theta(u_j,\ldots,u_N),\qquad
 p_j=P_\theta\left(\inf_{\ell\geq j}u_\ell\right).
\]
\Needspace{8\baselineskip}
All these envelopes lie between $\varphi-C$ and $\varphi+C$. The contact theorem and ordered contact give a finite-rooftop inequality: partition the contact set according to the first index $\ell$ for which $p_{j,N}=u_\ell$. On that piece the measure of $p_{j,N}$ is bounded by that of $u_\ell$, and $\mathrm{e}^{p_{j,N}}=\mathrm{e}^{u_\ell}$. Therefore
\[
 \int_X|\mathrm{e}^{p_{j,N}}-\mathrm{e}^u|\,
            \mathrm{MA}^{\mathrm{ext}}_\theta(p_{j,N})
 \leq\sum_{\ell=j}^N I_\ell\leq2^{1-j}.
\]
This uses only positive measures, including when their total masses are infinite.

As $N\to\infty$, $p_{j,N}\downarrow p_j$; the limit is the envelope of the infimum because every competitor is below every finite rooftop. The sequence $p_j$ increases to a genuine canonical limit $p$. Both passages hold in capacity by \cref{lem:monotone-psh-capacity}. Apply \cref{lem:weighted-extended-lsc} to the bounded nonnegative weights $|\mathrm{e}^{p_{j,N}}-\mathrm{e}^u|$, first in $N$ and then in $j$. We obtain
\[
 \int_X|\mathrm{e}^p-\mathrm{e}^u|\,\mathrm{MA}^{\mathrm{ext}}_\theta(p)=0.
\]
The inequalities $p_j\leq u_\ell$ for all $\ell\geq j$ imply $p\leq u$, first almost everywhere and then for canonical representatives. Also $p\sim u\sim\varphi$. Off their poles, $\mathrm{e}^p<\mathrm{e}^u$ exactly on $\{p<u\}$, so
$\mathrm{MA}^{\mathrm{ext}}_\theta(p)(\{p<u\})=0$.
The competitor $P_{\delta\theta}(\varphi)-C$ makes $P_{\delta\theta}(p)$ genuine. Apply \cref{thm:singular-strict-domination} with coefficient zero to obtain $p\geq u$, and hence $p=u$.

Finally $U_j=(\sup_{\ell\geq j}u_\ell)^*$ decreases to $u$ by Hartogs' lemma and the almost-everywhere extraction. The inequalities
$p_j\leq u_j\leq U_j$, together with the monotone capacity convergence of both outer sequences, prove the claim.
\end{proof}

\Needspace{10\baselineskip}
\begin{proposition}\label{prop:lower-barrier-criterion}
In \cref{thm:common-barrier-criterion}, one may replace the two bounds on $u_j$ by
\[
 \varphi-C\leq u_j\leq C
\]
provided the $L^1$ limit satisfies $u\preceq\varphi$. In particular the approximants may be bounded while $\varphi$ is unbounded.
\end{proposition}
\begin{proof}
The $L^1$ limit satisfies $u\geq\varphi-C$, so the additional hypothesis implies $u\sim\varphi$. Repeat the finite rooftop construction in the theorem. Every $p_{j,N}$ remains genuine and lies between $\varphi-C$ and $C$, which is enough for contact support and for the bounded exponential weights. The resulting inequality is still
\[
 \int_X|\mathrm{e}^{p_{j,N}}-\mathrm{e}^u|\,
             \mathrm{MA}^{\mathrm{ext}}_\theta(p_{j,N})\leq2^{1-j}.
\]
The limits $p_j$ exist by the common lower bound, and $p_j\leq u$ follows from $p_j\leq u_\ell$ for every $\ell\geq j$ and almost-everywhere convergence. Thus
$\varphi-C\leq p_j\leq u$, and all $p_j$ and their increasing limit $p$ have the same ordinary type as $\varphi$. Monotone capacity convergence and weighted lower semicontinuity give
$\mathrm{MA}^{\mathrm{ext}}_\theta(p)(\{p<u\})=0$. Singular domination applies to $(p,u)$ exactly as before, so $p=u$. The upper tail envelopes decrease to $u$, and the same sandwich gives capacity convergence. No singularity-type assertion was needed for the finite rooftops or for the individual approximants.
\end{proof}

\Needspace{12\baselineskip}
\begin{corollary}\label{cor:common-barrier-continuity}
Under the geometric and common-barrier hypotheses of \cref{thm:common-barrier-criterion}, suppose either
\begin{enumerate}
\item $\mathrm{MA}^{\mathrm{ext}}_\theta(u_j)\leq\nu$ for a fixed finite non-pluripolar measure $\nu$; or
\item $\mathrm{MA}^{\mathrm{ext}}_\theta(u_j)\leq A\omega_{w_j}^n$, where $A$ is fixed and $w_j\in\PSH(X,\omega)$ are uniformly bounded and converge in capacity to a bounded $w$.
\end{enumerate}
Then $u_j\to u$ in capacity, all the measures have finite mass, and
\[
 \mathrm{MA}^{\mathrm{ext}}_\theta(u_j)
 \rightharpoonup\mathrm{MA}^{\mathrm{ext}}_\theta(u).
\]
\end{corollary}
\begin{proof}
In the first case \cref{lem:fixed-npp-integrals} applied to exponentials proves \eqref{eq:common-barrier-self-error}. In the second case use \cref{lem:moving-bounded-integrals}, since the exponentials are bounded and have a common lower curvature bound. The criterion gives capacity convergence in both cases.

The masses are uniformly bounded by the corresponding dominating measures. For the deep tails use $u_j\geq\varphi-C$. In the first case their masses are at most $\nu(\{\varphi<C-T\})$, which tends to zero because $\nu$ ignores the pole set of $\varphi$. In the second case \cref{lem:bounded-mixed-capacity} bounds them by
$C'\operatorname{Cap}_\omega(\{\varphi<C-T\})$, which also tends to zero. Apply the raw-tail continuity theorem \cref{thm:undamped-tail-continuity} in its general-background form. It gives finiteness and weak convergence for the limiting measure as well.
\end{proof}

Its reference need not be an envelope fixed point; the smaller-background condition is enough. A model reference becomes useful when an a priori estimate supplies the two uniform barriers. If such a lower barrier is unavailable, constructing one in a relative envelope class is an additional problem, treated in \cref{sec:general-measure-relative}.
